% secciones/discusion.tex -- Sec 4 Discusión

Este estudio parte de una pregunta simple ---cuánto del CO$_2$ medido en un contenedor de bioconversión respiran las larvas y cuánto el lecho--- y la aborda mediante una implementación verificada del núcleo fisiológico, una comparación entre configuraciones y una evaluación con límites declarados. El resultado central es que C mejora el ajuste del caso de estudio, pero la evaluación cruzada no demuestra superioridad predictiva general. Su RMSE censurado agregado es 130.97 frente a 180.53~ppm/min de E0; sin embargo, los parámetros de prepupa permanecieron fijos en valores ajustados con todos los datos tardíos, por lo que el agregado es híbrido. En los cinco cierres tempranos, A presenta el menor RMSE censurado, 122.87~ppm/min, frente a 147.06 de C; en los tres tardíos, C alcanza 98.50 frente a 168.00~ppm/min de E0, pero esos cierres participaron en la calibración global de prepupa. R reduce el error agregado a 164.41~ppm/min, mientras que el paso de R a C combina lecho, acople, pastoreo y prepupa y no aísla una contribución mecanística. Además, el intervalo percentil bootstrap de la reducción agregada de C frente a E0, $(-42.8,46.0)$~\%, incluye valores negativos y hereda la dependencia tardía; no constituye una prueba concluyente de superioridad.

El contraste estructural es compatible con un control top-down de la biomasa microbiana dentro de la formulación asumida. En calibración, anular el pastoreo en S eleva el RMSE censurado a 435.31~ppm/min frente a 86.41~ppm/min de C, conservando el resto de la estructura. En la evaluación cruzada, A y C difieren además en la recalibración del núcleo y en el acople con el alimento, y el orden depende del periodo: A presenta el menor RMSE censurado temprano, mientras que el menor valor tardío de C está condicionado por la calibración global de prepupa. Por tanto, esa comparación tampoco aísla el efecto del pastoreo. No se descartan mecanismos alternativos de decaimiento endógeno capaces de reconciliar controles y tratamientos, y la posición interior de $\gamma=0.5798$ no basta para establecer identificabilidad. La interpretación biológica propuesta para el término es la ingesta de microorganismos junto con el sustrato; su destino de carbono y una posible fracción asimilada requieren verificación experimental.

La evaluación externa muestra un desempeño dependiente del sistema y del nivel de ajuste. En \textcite{ermolaevGreenhouseGasEmissions2019}, E0 reproduce el acumulado y los flujos dentro de los criterios declarados en transferencia pura. En ese mismo nivel A, C mantiene el acumulado dentro del margen ($+17.9$~\%), pero su RMSE de flujos es 1.66~mg/min (30.8~\%), fuera del criterio; su nivel B constituye una interpolación con $n=p$ y mejora el gas a costa de subestimar la biomasa final. En \textcite{schonUsingGasExchange2026}, C queda dentro del margen declarado en tres de cuatro dietas después de reajustar $f$ para cada dieta; la dieta más pobre permanece fuera del margen. En consecuencia, estos resultados respaldan una transferibilidad condicionada al sistema y a la calibración local, no una validez universal del mecanismo ni del conjunto paramétrico.

Los límites del estudio delimitan la interpretación: ocho cierres de tratamiento, cuatro de ellos censurados; ausencia de réplicas biológicas de gas; cinco estimaciones en fronteras de búsqueda; e independencia limitada de la fase tardía. El crecimiento tardío no entró en la etapa~2, pero los cierres tardíos de gas sí participaron en ella, y los tres cierres exactos tardíos intervinieron en la calibración global de prepupa. En este experimento, los siete controles sin larvas fueron necesarios para calibrar la fuente microbiana; su cobertura no incluye la dinámica inicial ni el día larvario 14 de D1, de modo que no constituyen una validación temporal independiente. El lecho sin control larvario condujo a sobreestimaciones tardías, una observación propia del caso de estudio que no debe extenderse a todos los reactores sin evaluación adicional. Las réplicas biológicas independientes y el registro dinámico de la humedad son necesarios para estrechar la incertidumbre y examinar la transferencia futura.
